\section*{Chapter \ref{ch:qm:stochastic-differential-equations} --- Stochastic Differential Equations}

\begin{solution}{exo:qm:stochastic-differential-equations:1}
$\ln 2/2 = 0.35$ years, 87 trading days.
\end{solution}

\begin{solution}{exo:qm:stochastic-differential-equations:2}
$\P(\ln S_1 < \ln S_0) = \Phi(-(0.05 - 0.02)/0.20) = \Phi(-0.15) = 44.0\%$: the median grows at $\mu -
\sigma^2/2 = 3\%$.
\end{solution}

\begin{solution}{exo:qm:stochastic-differential-equations:3}
$\mathcal Lf = \kappa(\bar x - x)f' + \tfrac12\sigma^2f^{\prime\prime}$; for $f(x) = x$, $\mathcal Lf = \kappa(\bar x - x)$, so
$\frac{d}{dt}\E[X_t] = \kappa(\bar x - \E[X_t])$ and $\E[X_t] = \bar x + (X_0 - \bar x)e^{-\kappa t}$.
\end{solution}

\begin{solution}{exo:qm:stochastic-differential-equations:4}
$\kappa(\bar x - y)p = \tfrac12\sigma^2p'$ gives a Gaussian density with variance $\sigma^2/2\kappa = 0.16/4 = 0.04$.
\end{solution}

\begin{solution}{exo:qm:stochastic-differential-equations:5}
Dynkin with $f(v) = v$ gives $\frac{d}{dt}\E[v_t] = \kappa(\bar v - \E[v_t])$, so $\E[v_{0.5}] = 0.04 + 0.05e^{-1} =
0.0584$.
\end{solution}

\begin{solution}{exo:qm:stochastic-differential-equations:6}
$B = (1 - e^{-2.5})/0.5 = 1.8358$; $A = (0.04 - 0.0002)(1.8358 - 5) - 0.0001 \times 1.8358^2/2 = -0.1261$;
$e^{A - 0.03B} = e^{-0.1812} = 0.8343$.
\end{solution}

\begin{solution}{exo:qm:stochastic-differential-equations:7}
24.5\% at a Feller ratio of one: the condition keeps the exact process away from zero, not the
Gaussian Euler step. In the table, $\eta = 0.25$ (ratio 2.56) gives 0.08\%, and $\eta = 0.3$ (ratio 1.78)
already 1.6\%.
\end{solution}

\begin{solution}{exo:qm:stochastic-differential-equations:8}
Flooring removes the NaN but not the error: every step that would have crossed zero is replaced by
zero, which then has zero diffusion, so the scheme adds mass at zero and biases the law near the
boundary. It is still a first-order Euler scheme, not the exact transition; check its moments
against $\E[v_t] = \bar v + (v_0 - \bar v)e^{-\kappa t}$ and the Gamma stationary law, or use the exact or
full-truncation scheme.
\end{solution}

\begin{solution}{pb:qm:stochastic-differential-equations:1}
\textbf{1.} $2 \times 2 \times 0.04/0.36 = 0.44$: below one, zero is attainable.
\textbf{2.} $\mathrm{Gamma}(0.44, 0.09)$: mean 0.04, standard deviation $\sqrt{0.04 \times 0.36/4} = 0.06$.
\textbf{3.} 27.9\%.
\textbf{4.} $\eta = \sqrt{2 \times 2 \times 0.04} = 0.4$.
\textbf{5.} $\ln 2/2 = 0.35$ years.
\textbf{6.} $\Phi(-(0.004 + 2 \times 0.036/252)/(0.6\sqrt{0.004/252})) = 3.6\%$.
\textbf{7.} 75.2\% of 40\,000 paths.
\textbf{8.} 75.3\% with four steps a day, 74.6\% with sixteen, 75.2\% with weekly steps.
\textbf{9.} The exact process comes arbitrarily close to zero on those paths; from a small $v$ a
Gaussian step of any length has a fixed chance of crossing, and more steps near zero means more
chances.
\textbf{10.} It propagates through every sum and average into the book's value, and a comparison
with a NaN is false, so limit checks can pass silently.
\textbf{11.} Minimum 0 (never negative); one-year mean 0.0394 against $\bar v = 0.04$.
\textbf{12.} Mean 0.0393; at zero on 5.2\% of the steps.
\textbf{13.} The exact scheme has no discretisation bias; full truncation does, and its bias
vanishes as $\Delta t \to 0$.
\textbf{14.} Reflection turns every overshoot into a positive value of the same size, pushing mass
up and biasing the mean upward near zero.
\textbf{15.} The exact sampler is slow and the Euler variants need small steps; the
quadratic-exponential scheme (One Quant Book~5, chapter~23) is accurate at daily or larger steps.
\textbf{16.} No: the market's smile asks for a high vol-of-vol, and imposing the condition would
distort the fit to save a numerical scheme; fix the scheme.
\textbf{17.} No NaN and no negative variance on any seed; $\E[v_t]$ and $\Var(v_t)$ against their closed
forms; the stationary law against the Gamma density; convergence as the step is refined.
\textbf{18.} The root cause (plain Euler with a Feller ratio of 0.44), the size of the failure (75\% of
paths), the fix (exact or full truncation) and its validation, and the tests added.
\textbf{19.} Named result: \emph{the NaN in the overnight run}: with a Feller ratio of 0.44, 75\% of
daily plain-Euler paths produce a negative variance within a year, whatever the step; the exact
and full-truncation schemes produce none.
\textbf{20.} It keeps the exact process away from zero; it does not keep a Gaussian step from
crossing it.
\end{solution}

\begin{solution}{iq:qm:stochastic-differential-equations:1}
$S_T = S_0\exp((\mu - \tfrac12\sigma^2)T + \sigma W_T)$; the median is $S_0e^{(\mu - \sigma^2/2)T}$, below the mean
$S_0e^{\mu T}$.
\iqlookfor{Itô on $\ln S$ and the lognormal law.}
\end{solution}

\begin{solution}{iq:qm:stochastic-differential-equations:2}
$\kappa = \ln 2/10 = 0.069$ a day; after 30 days, three half-lives, $2^{-3} = 12.5\%$ remains.
\iqlookfor{$e^{-\kappa t}$ and half-lives.}
\end{solution}

\begin{solution}{iq:qm:stochastic-differential-equations:3}
$2\kappa\bar v \ge \eta^2$: then the drift near zero beats the diffusion and zero is never reached;
otherwise the process touches zero and reflects. The law stays nonnegative either way, but an Euler
step from near zero can go negative, and more so when the condition fails.
\iqlookfor{the boundary classification and its numerical consequence.}
\end{solution}

\begin{solution}{iq:qm:stochastic-differential-equations:4}
$u(t,x) = \E[e^{-\int_t^Tr}g(X_T) \mid X_t = x]$ solves $\partial_tu + \mathcal Lu - ru = 0$, $u(T) = g$. For
$dS = \mu S\,dt + \sigma S\,dW$ and constant $r$: $\partial_tu + \mu S\partial_Su + \tfrac12\sigma^2S^2\partial_{SS}u - ru = 0$.
\iqlookfor{generator plus discounting, terminal condition.}
\end{solution}

\begin{solution}{iq:qm:stochastic-differential-equations:5}
Make the failure reproducible (log the seed and the path index), find the first non-finite value
and the step that produced it, and read the scheme at that state. Then make it impossible: exact or
boundary-respecting schemes, assertions on the state space inside the kernel, a test that runs many
seeds and fails on any NaN, and aggregation that refuses non-finite inputs instead of propagating
them.
\iqlookfor{reproducibility, root cause, and invariants enforced in code.}
\end{solution}

\begin{solution}{iq:qm:stochastic-differential-equations:6}
A \omterm{def:qm:stochastic-differential-equations:sde}{strong solution} is built on a given \omterm{def:qm:brownian-motion:bm}{Brownian motion}; a weak one only asks for some probability
space carrying a process and a \omterm{def:qm:brownian-motion:bm}{Brownian motion} that satisfy the equation. Tanaka's equation $dX =
\operatorname{sign}(X)\,dW$, $X_0 = 0$, has \omterm{def:qm:stochastic-differential-equations:sde}{weak solutions} (any \omterm{def:qm:brownian-motion:bm}{Brownian motion} $X$ with $W = \int\operatorname{sign}(X)\,dX$)
but no strong one, because $X$ cannot be recovered from $W$.
\iqlookfor{the definitions and Tanaka's example.}
\end{solution}
